\documentclass[a4paper,11pt]{article}
\usepackage[utf8]{inputenc}
\usepackage[T1]{fontenc}
\usepackage[french]{babel}
\usepackage{lmodern}
\usepackage{amsmath,amssymb,amsthm}
\usepackage{mathrsfs}
\usepackage{enumitem}
\usepackage{multicol}

\setlist[enumerate]{label=\textbf{\textcolor{Couleur8}{\arabic*.}}, left=1em}

% Si vous voulez aussi modifier les listes enumerate imbriquées :
\setlist[enumerate,2]{label=\textcolor{Couleur8}{\alph*)}}
\setlist[enumerate,3]{label=\textcolor{Couleur8}{\roman*)}}
\setlist[enumerate,4]{label=\textcolor{Couleur8}{\Alph*)}}
\usepackage{graphicx}
\usepackage{xcolor}
\usepackage{tcolorbox}
\usepackage{geometry}
\usepackage{fancyhdr}
\usepackage{titlesec}
\usepackage{cancel}
\usepackage{ifthen}
\usepackage{tikz}
\usetikzlibrary{arrows.meta}
\usepackage{tkz-tab}
\usepackage{eurosym}

% OPTION PROF/ÉLÈVE
% Décommentez UNE des deux lignes suivantes :
\newboolean{prof}\setboolean{prof}{true}   % Version professeur (avec corrections des devoirs)
%\newboolean{prof}\setboolean{prof}{false}  % Version élève (sans corrections des devoirs)

\definecolor{Couleur1}{HTML}{4A5D7A} %Th
\definecolor{Couleur2}{HTML}{A5B5D6} %Prop
\definecolor{Couleur3}{HTML}{A8C68A} %Méthode
\definecolor{Couleur6}{HTML}{8A9CC1} %Def
\definecolor{Couleur7}{HTML}{E6B459} %Rq Voc
\definecolor{Couleur8}{HTML}{D36740} %Titres
\definecolor{correction}{HTML}{D36740} % Couleur pour les corrections
\definecolor{coopmathscorr}{HTML}{F15929} % Couleur corrections coopmaths

% Configuration des marges
\geometry{
	top=2cm,
	bottom=2cm,
	inner=1.5cm,
	outer=1.5cm,
}

% Police
\renewcommand{\familydefault}{\sfdefault}

% Compteurs
\newcounter{seancenum}
\newcounter{seancenumD}
\setcounter{seancenum}{1}
\setcounter{seancenumD}{1}

% Configuration des en-têtes et pieds de page
\pagestyle{fancy}
\fancyhf{}
\ifthenelse{\boolean{prof}}{
    \fancyhead[L]{\textcolor{Couleur8}{\textbf{Corrections Automatismes - Classe de seconde - Version Professeur}}}
}{
    \fancyhead[L]{\textcolor{Couleur8}{\textbf{Corrections Automatismes - Séance 30 - Classe de seconde}}}
}
\fancyhead[R]{\textcolor{Couleur8}{\textbf{2025-2026}}}
\fancyfoot[L]{\textcolor{Couleur8}{Mme Bertail et Mme Pinto}}
\fancyfoot[R]{\textcolor{Couleur8}{\thepage}}
\renewcommand{\headrulewidth}{1pt}
\renewcommand{\headrule}{\hbox to\headwidth{\color{Couleur8}\leaders\hrule height \headrulewidth\hfill}}

% Environnements personnalisés
\newtcolorbox{seance}[1]{
    colback=Couleur6!10,
    colframe=Couleur1!65,
    fonttitle=\bfseries\large,
    title=Correction Séance \theseancenum #1,
    left=5mm,
    right=5mm,
    top=3mm,
    bottom=3mm,
    arc=0mm,
    after=\stepcounter{seancenum}
}

\newtcolorbox{seanceD}[1]{
    colback=Couleur6!5,
    colframe=Couleur6!45,
    fonttitle=\bfseries\large,
    title=Correction Devoirs - Séance \theseancenumD #1,
    left=5mm,
    right=5mm,
    top=3mm,
    bottom=3mm,
    arc=0mm,
    after=\stepcounter{seancenumD}
}

\begin{document}

\setcounter{seancenum}{30}
\newpage
\begin{seance}{} %30

\begin{enumerate}
\item $\dfrac{x}{3}-2x=\dfrac{x}{3}-\dfrac{2x\times 3}{3}=\dfrac{x}{3}-\dfrac{6x}{3}=\dfrac{x-6x}{3}=\dfrac{-5x}{3}=$ ${\color[HTML]{F15929}\boldsymbol{\dfrac{-5}{3}x}}$\\La bonne réponse est la réponse {\bfseries \color[HTML]{F15929}D}.

\item On supprime les parenthèses en changeant les signes dans la deuxième parenthèse (précédée du signe $-$) :\\$A= -5z+2 -8z^2-7z-5$\\On réduit :\\$A={\color[HTML]{F15929}\boldsymbol{-8z^2-12z-3}}$\\La bonne réponse est la réponse {\bfseries \color[HTML]{F15929}D}.

\item La fonction $g$ s'annule en $-3$ et la fonction $h$ s'annule en $1$.\\
    Quand la droite est en-dessous de l'axe des abscisses, la fonction est négative et quand elle est au-dessus, la fonction est positive.\\
    On en déduit le tableau de signes de leur produit :\\
    \begin{center}
    \begin{tikzpicture}[baseline, scale=0.5]
    \tkzTabInit[lgt=5,deltacl=1,espcl=3]{ $x$ / 2, $g(x)$ / 2, $h(x)$ / 2, $f(x)$ / 2}{ $-\infty$, $-3$, $1$, $+\infty$}
	\tkzTabLine{  , +, z, -, t, -}
	\tkzTabLine{  , -, t, -, z, +}
	\tkzTabLine{  , -, z, +, z, -}
	\end{tikzpicture}
    \end{center}
La bonne réponse est la réponse {\bfseries \color[HTML]{F15929}B}.

\item Après avoir tiré un jeton rouge, il reste $8$ jetons rouges et $8$ jetons noirs, soit $16$ jetons au total.\\La probabilité que le deuxième jeton soit noir sachant que le premier est rouge est : $\dfrac{8}{16}={\color[HTML]{F15929}\boldsymbol{\dfrac{1}{2}}}$.\\La bonne réponse est la réponse {\bfseries \color[HTML]{F15929}D}.

\item L'écart interquartile est la différence entre le troisième quartile et le premier quartile.\\
      D'après le diagramme en boite, on a $Q_3 = 55$ et $Q_1 = 20$.\\
      Donc l'écart interquartile est $Q_3 - Q_1 = 55 - 20 = {\color[HTML]{F15929}\boldsymbol{35}}$.\\La bonne réponse est la réponse {\bfseries \color[HTML]{F15929}C}.

\item On observe sur le graphique que le coefficient directeur est positif ($D$ représente une fonction affine croissante) et que l'ordonnée à l'origine est strictement positive ($D$ coupe l'axe des ordonnées au-dessus de l'origine).\\On écrit les équations qui ne sont pas forme réduite, sous forme réduite :\\$\bullet\:$ $y=x^2-(x-3)^2+12$ s'écrit $y=3x+3$.\\$\bullet\:$ $y=-3x+3$ est  sous forme réduite.\\$\bullet\:$ $6x+2y-6=0$ s'écrit $y=-3x+3$.\\$\bullet\:$ $y=3x-3$ est  sous forme réduite.

\medskip
La seule équation ayant un coefficient directeur positif et une ordonnée à l'origine positive est : ${\color[HTML]{F15929}\boldsymbol{y=x^2-(x-3)^2+12}}$.\\La bonne réponse est la réponse {\bfseries \color[HTML]{F15929}C}.

\end{enumerate}

\end{seance}

\end{document}